% Generated draft: requires mathematical and source review.
\documentclass[11pt]{article}
\usepackage[margin=25mm]{geometry}
\usepackage{fontspec}
\setmainfont{Latin Modern Roman}
\usepackage{amsmath,amssymb,mathtools,unicode-math}
\setmathfont{Latin Modern Math}
\usepackage{enumitem,needspace,xurl,hyperref}
\setlength{\parindent}{0pt}
\setlength{\parskip}{5pt}
\setlength{\emergencystretch}{3em}
\newcommand{\N}{\mathbb{N}}
\newcommand{\Z}{\mathbb{Z}}
\newcommand{\Q}{\mathbb{Q}}
\newcommand{\R}{\mathbb{R}}
\newcommand{\C}{\mathbb{C}}
\newcommand{\F}{\mathbb{F}}
\newcommand{\A}{\mathbb{A}}
\newcommand{\PP}{\mathbb{P}}
\newcommand{\E}{\mathbb{E}}
\newcommand{\eps}{\varepsilon}
\newcommand{\dd}{\,\mathrm{d}}
\begin{document}
\section*{UnsolvedMath ID: 20003264}
\textbf{Source problem code:} catalogue entry
\subsection{Definitions and mathematical statement}
This record is connective prose, not an autonomous conjecture. The supplied source inspection says it ends at the colon immediately before separately numbered Conjecture 3 in the 2004 AIM report. It first comments on a preceding Conjecture 2, then surveys local rigidity, then introduces global Anosov rigidity. Neither full conjecture is included, so their missing statements are not reconstructed here.

The source reports that [97] proves the preceding Conjecture 2 when $A$ is the full split Cartan subgroup, and attributes to A. Katok the suggestion to use the method of [39]. It describes the higher-rank semisimple-group/lattice local-rigidity program of R. Zimmer [209], mentioning Anosov-action results [85,94,95,97], weaker-hyperbolicity results [131], and results without hyperbolicity [60]. Its next paragraph motivates a folklore conjecture using infranilmanifold Anosov examples and topological conjugacy to hyperbolic automorphisms [62,127], then cites [130] for the following conjecture. These are historical context and references with no bibliography supplied here. The phrase about the ``only known examples'' must not be promoted to an updated literal classification of smooth diffeomorphisms; the currency symbol in its ``100,000'' phrase is extraction-damaged.

The aid develops a separate background distinction between topological and $C^1$ conjugacy. Let $A\in GL(d,\Z)$ be hyperbolic, meaning it has no eigenvalue of modulus one, and let $L_A:T^d\to T^d$ be the induced toral automorphism. For all sufficiently small nonzero real $\eps$, choose a smooth diffeomorphism $\psi_\eps$ supported in a coordinate ball at $0$, with
\[
 \psi_\eps(0)=0,\qquad D\psi_\eps|_0=(1+\eps)I,
 \qquad\psi_\eps\longrightarrow\mathrm{id}\text{ in }C^1.
\]
For example a compactly supported coordinate bump with value one near zero provides this family. Put $f_\eps=\psi_\eps\circ L_A$. The aid invokes openness and structural stability of Anosov diffeomorphisms: $f_\eps$ is smooth Anosov, arbitrarily $C^1$-close to $L_A$, and topologically conjugate to it.

It nevertheless has no conjugacy to $L_A$ by a $C^1$ diffeomorphism. Any such $h$ would carry its fixed origin to a fixed point of $L_A$ and give the similarity relation
\[
 Dh|_0\,Df_\eps|_0=A\,Dh|_0.
\]
But $Df_\eps|_0=(1+\eps)A$, so its determinant is $(1+\eps)^d\det A\ne\det A$ for the chosen small nonzero $\eps$. This violates similarity. The obstruction requires the conjugating map to be a $C^1$ diffeomorphism with invertible derivative, not merely a possibly singular differentiable homeomorphism.

This single-diffeomorphism perturbation is contextual mathematics. It neither supplies the missing folklore conjecture nor refutes the higher-rank action-rigidity assertions, which impose additional action relations and hypotheses.
\subsection{Short English statement}
The extracted entry joins historical rigidity discussion and stops before the actual next conjecture. Its attached aid illustrates why topological conjugacy is weaker than $C^1$ diffeomorphic conjugacy: an arbitrarily small Anosov perturbation of a toral automorphism changes a fixed-point derivative determinant, obstructing a $C^1$ conjugacy while structural stability preserves a topological one.
\subsection{Sources}
\begin{itemize}
\item Local UnsolvedMath registry, permanent ID $20003264$, source code catalogue entry; canonical location \texttt{aim-other-notes.json}, \texttt{notes[70]}.
\item AIM workshop report \emph{Emerging Applications of Measure Rigidity}, 2004, \url{https://aimath.org/WWN/measrigid/measrigid.pdf}. The supplied source assessment identifies this as connective prose preceding Conjecture 3, despite the local conjecture/open tags and reconstructed/unverified label.
\item Source labels [97] (full split-Cartan case), [39] (method suggested by A. Katok), [209] (R. Zimmer program), [85,94,95,97] (Anosov actions), [131] (weaker hyperbolicity), [60] (without hyperbolicity), [62,127] (infranilmanifold conjugacy), and [130] (following folklore conjecture). No full entries are supplied, so none is guessed.
\item Research attempt 1 is marked context-only/not-a-problem, proved, and low novelty confidence, with an explicit machine-generated, non-peer-reviewed research-aid caveat requiring independent verification before citation. Structural stability and the derivative obstruction are its developed background, not the missing source question.
\end{itemize}
\end{document}
